\appendix
\section{Additional tables, provenance, and corrections}
\label{sec:appendix}

This appendix carries the full arm matrix of \S\ref{sec:eval-crosslang}, the data provenance of
\S\ref{sec:eval}, and the correction log of the instrumentation faults of
\S\ref{sec:eval-correct}.  \S\ref{sec:appendix:limits} then states, claim by claim, what the Lean
development proves and what it does not.

\subsection{The full arm matrix}
\label{sec:appendix:arms}

\tab\ref{tab:allarms} reports every arm of the cross-language matrix, in the same units as
\tab\ref{tab:crosslang} and from the same run.  The four shapes appear in descending order of
allocation intensity, and within a shape the arms are ordered by language and configuration.  Some
entries repay reading individually: the leaky C++ arm on \code{promote} runs in 109\,\ms and holds
143 MB, which is what a bump allocator with no reuse actually costs; the Go arms are competitive on
time on two shapes but pay through the garbage collector in the others; and Node on \code{chase} is
the one entry three orders of magnitude off, which is what a pointer-chasing loop looks like when
every object is a boxed heap record.  The ratios of \tab\ref{tab:crosslang} are formed against the
fastest arm of a comparison language, not against an average arm, so they understate the advantage
against a typical configuration and overstate nothing.

\begin{table*}[t]
\centering
\setlength{\tabcolsep}{4pt}
\scriptsize
\caption{All eleven arms: milliseconds (minimum / median) and peak RSS in MB, per shape.}
\label{tab:allarms}
\resizebox{\ifdim\width>\linewidth\linewidth\else\width\fi}{!}{
\begin{tabular}{@{}lrrrr@{}}
\toprule
arm & promote & sizeratio & chase & cycleratio \\
\midrule
extc\texttt{:}default & 39.0 / 40.4 / 7.5 & 16.0 / 16.6 / 7.5 & 117.3 / 129.2 / 66.2 & 15.0 / 15.3 / 7.4 \\
extc\texttt{:}off & 82.8 / 84.1 / 107.0 & 59.5 / 61.3 / 154.2 & 104.0 / 107.5 / 60.7 & 21.0 / 22.2 / 43.4 \\
cpp\texttt{:}default & 109.4 / 113.3 / 143.1 & 77.7 / 79.6 / 187.0 & 119.2 / 122.8 / 79.6 & 56.0 / 57.7 / 115.3 \\
cpp\texttt{:}careful & 109.9 / 111.8 / 143.1 & 17.9 / 18.8 / 7.5 & 116.0 / 121.2 / 75.8 & 36.1 / 36.4 / 7.5 \\
rust\texttt{:}default & 39.1 / 40.2 / 7.4 & 20.6 / 21.3 / 7.5 & 220.2 / 227.3 / 90.0 & 49.8 / 51.2 / 113.4 \\
rust\texttt{:}careful & 39.5 / 40.3 / 7.5 & 20.6 / 21.4 / 7.5 & 219.3 / 224.9 / 90.2 & 17.9 / 18.7 / 7.5 \\
go\texttt{:}default & 106.1 / 114.1 / 7.5 & 45.6 / 46.0 / 7.5 & 146.3 / 154.7 / 64.6 & 33.7 / 35.0 / 7.6 \\
go\texttt{:}off & 128.8 / 132.1 / 112.6 & 80.8 / 82.6 / 160.6 & 115.3 / 127.7 / 63.6 & 38.9 / 39.5 / 46.0 \\
java\texttt{:}default & 103.0 / 108.4 / 127.7 & 94.1 / 99.6 / 127.6 & 174.1 / 191.6 / 114.6 & 69.1 / 78.4 / 104.3 \\
node\texttt{:}split & 102.8 / 108.1 / 74.7 & 85.3 / 91.4 / 72.0 & 1119.6 / 1179.4 / 355.8 & 60.6 / 61.9 / 70.5 \\
node\texttt{:}bigint & 166.7 / 170.4 / 82.9 & 113.2 / 117.4 / 79.6 & 1321.1 / 1361.8 / 353.1 & 90.6 / 92.3 / 70.8 \\
\bottomrule
\end{tabular}}
\end{table*}


\subsection{Provenance and reproduction}
\label{sec:appendix:artifact}

\label{sec:artifact}

\paragraph{What is in the repository.}
The compiler is built with \code{make} into \code{build/extc}; it needs a C compiler and nothing
else, and the generated C links only against libc.  The runtime is not a library: \code{src/back/reclaim.c}
is included into every translation unit that asks for it, and a unit that the analysis judges idle
takes none of it.  The four shape programs and the eleven comparison arms are in
\code{bench/gc/v2}, the pause harness in \code{bench/gc/v2/pause}, the driver test suite in
\code{tests/reclaim}, the emitted-code gates in \code{tools}, and the Lean development in
\code{proofs}.

\paragraph{How to re-run each table.}
The harness is one shell script plus one report script, and every table in the paper is a rendering
of the \code{.tsv} that the run writes:
\begin{itemize}
\item \tab\ref{tab:crosslang} and \tab\ref{tab:allarms}: \code{run.sh align} builds every arm
into the directory named by \code{B}, then \code{run.sh bench} with \code{REPS=7} writes the matrix
that \code{report.py} renders.
\item \tab\ref{tab:pause}: \code{pause/extc.sh} with \code{REPS=2} writes the pause logs, and
\code{pause_report.py} renders \code{PAUSE.md} and \code{pause\_v2.tsv}.
\item \tab\ref{tab:ablate} and \tab\ref{tab:ablatepause}: \code{measure\_ablation.py}, with
\code{--pause} for the second.  The script generates the C for each shape once, then recompiles it
with one \code{-D} constant changed and checks that every variant prints the same bytes as the
reference before timing it.
\item \tab\ref{tab:bounded}: \code{measure\_bounded.py} runs \code{bounded.extc} at three orders of
magnitude of round count, with reclamation on and off.
\item \tab\ref{tab:census}: \code{census.py} compiles every \code{.extc} under \code{tests/} one
file at a time and reads the tier of the emitted header.
\item the two figures: \code{gen\_tables.py} and \code{gen\_roofline.py}, both reading only
\code{paper/data/*.tsv}.
\end{itemize}


\paragraph{Frozen build.}
Every number in this paper is from the tree at commit \code{98bd56d2}, in which the runtime is
\code{src/back/reclaim.c} (4,212 lines, md5 \code{7151a431}), the code generator is
\code{src/back/codegen.c} (16,017 lines, md5 \code{2f560e2e}), and the driver suite is
\code{tests/reclaim/driver.c} (1,637 lines, md5 \code{28dd4009}).  The driver reports 216 passing
cases, and the whole suite is also run with \code{-fsanitize=address}, with
\code{-fsanitize=undefined} and with the runtime forced out of line, so a correctness result is not
an artifact of one build configuration.  Two properties are held by gates rather than by tests: the
instruction count of the generated code for each shape (a frozen budget in the whole-program run)
and the byte-identity of the emitted C for every shape, so that a change to an unrelated section of
the compiler cannot silently move the measurements.

\paragraph{Provenance of the numbers.}
\code{bench/gc/v2/matrix.tsv} is the cross-language run and \code{bench/gc/v2/pause\_v2.tsv} the
pause run, both checked in.  \code{paper/data/} holds copies of these and the three files produced
by the mechanism-ablation, boundedness and census scripts, together with the model output of
\code{gen\_roofline.py}.  The tables in \code{paper/tables} are generated from those files and are
not edited by hand; where a table and a \code{.tsv} disagree the \code{.tsv} is correct.  The
numbers that are transcribed rather than generated are the constants of
\tab\ref{tab:calib}, which come from the calibration runs recorded in \code{../extc-docs/devlog/analysis}, and
\code{paper/data/roofline-pareto.tsv} which is model output rather than measurement.


\subsection{Corrections to earlier internal measurements}
\label{sec:appendix:corrections}

Progress in this project was measured by instruments we wrote, so a fair paper has to say which
instruments were later found wrong and what that does to the record.  Three faults were found, all
in the pause and drain probes, and each one is fixed with a gate that fails if it recurs.

\begin{enumerate}
\item \emph{The histogram bucket that was always true.}  The distribution probe classified each
sample by comparing cumulative bounds but the last bucket was assigned rather than compared, so the
condition "over one millisecond" held for every sample.  The fault cannot make a bound look better
than it is, but it destroys the resolution needed to see a step-landing improvement, and it was
present in the internal report where a 60\,\us target appeared to be met.  The probe now compares and
\tab\ref{tab:pause} is produced by it.
\item \emph{The clock that was silently zero.}  The step timer asked for per-thread CPU time without
declaring the feature-test macro that exposes it on this platform, so the clock call failed and the
probe reported zero for every step.  A zero population is indistinguishable from a perfect one, which
is why the gate now treats a missing field as a failure rather than a pass, and why the pause section
reports wall time by default and per-thread CPU time only as a diagnostic.
\item \emph{The queue that was assumed to drain in one pass.}  A driver test asserted that a single
bounded drain leaves the deferred-free queue empty.  The budget is eight slabs or one hundred
\us, whichever comes first, so under sanitizer load the pass sometimes emptied one of two queued
slabs and the assertion failed even though the runtime was correct.  The judgement was not weakened:
the test now triggers drain at several allocation points until the queue is empty and still requires
the final state to be exactly zero, and the full drain path is asserted separately.
\end{enumerate}

The general lesson is in \S\ref{sec:eval-correct}: a measurement that is cheap to take is also cheap
to break in a way that looks like good news, so the gates are written to fail on absence of evidence
as well as on contradictory evidence.
